\chapter{Device Authentication Hardening}
\label{ch:4.2}
\chaptertopic{Topic 4.2}

\chapterquote{Treat your password like your toothbrush. Don't let anybody else use it, and get a new one every six months.}{commonly attributed to Clifford Stoll}

\section{Motivation}

Chapter~\ref{ch:4.1} covered how an endpoint gets compromised. This chapter covers the layer that
decides who is allowed to touch it in the first place, and what happens after they successfully do.
Four topics build on each other in sequence. First, the password itself. Second, the factor that
backs it up when it inevitably leaks. Third, the session that keeps a login active without
re-asking for a password on every click. Fourth, the access-control model that decides how much a
successfully authenticated account is actually allowed to do. Each layer assumes the one before it
can fail.

\begin{apalignment}{Unit 4: Securing Devices}{Topic 4.2 -- Authentication}
\begin{apobjectives}
\apobjective{4.2.A}{Analyze Risk}{Account for why a system keeps the output of a one-way function
-- a hash, digest, or checksum -- in place of the password itself, and what that substitution does
and does not buy when the credential store is stolen.}
\apobjective{4.2.B}{Analyze Risk}{Explain the mechanism behind each family of attack on
credentials -- exhaustive search, wordlists, precomputed tables, reused breach data, and low-and-slow
guessing across many accounts -- identifying which weakness each one feeds on and which single
defense removes it.}
\apobjective{4.2.C}{Mitigate Risk}{Given a description of how a person proved who they were,
determine which categories of evidence were actually used, whether the combination constitutes more
than one independent factor, and how resistant the resulting scheme is to interception.}
\apobjective{4.2.D}{Mitigate Risk}{Distinguish session fixation from session hijacking by when in
the login lifecycle each attack plants or steals its token, and identify which session control --
regeneration, the Secure/HttpOnly/SameSite cookie flags, or timeout policy -- defeats which one.}
\apobjective{4.2.E}{Mitigate Risk}{Choose an access-control model (DAC, MAC, RBAC, or ABAC)
appropriate to a given scenario, and identify the organizational process that catches privilege
that accumulates even after an account was originally provisioned correctly.}
\end{apobjectives}
\end{apalignment}

\section{Passwords and Password Policy}

A password is the oldest authentication mechanism there is: a shared secret known only to the user
and the system. Two properties dominate its strength -- length and unpredictability -- and they do
not trade off equally. A \textbf{brute-force attack} tries every possible combination; a
\textbf{dictionary attack} tries words, names, and common substitutions. Both become exponentially
more expensive as length increases. Adding one character to a ten-character password drawn from a
95-character alphabet multiplies the search space by 95; adding four characters multiplies it by
roughly 81 million. Complexity rules -- mandatory uppercase, digits, symbols -- were introduced to
push users off dictionary words, but they are not a substitute for length. A long random
passphrase beats a short, symbol-sprinkled word against every attack class, because a predictable
substitution pattern (\texttt{p@ssw0rd}) generates almost no additional entropy.

\begin{definition}{Password entropy}
The number of bits of uncertainty in a password, computed as $H = L \times \log_2(R)$, where $L$ is
the password's length in characters and $R$ is the size of the character set (\emph{alphabet}) each
character is independently drawn from. Entropy measures search-space size, not memorability -- a
password can be long and still low-entropy if $R$ is small or the characters are not independently
chosen.
\end{definition}

\begin{example}[Computing entropy for two eight-character passwords]
Take \texttt{Tr0ub4dor\&3}-style complexity rules versus a long random passphrase, both compared at
equal length first. An eight-character password drawn from a 95-character keyboard alphabet
(lowercase, uppercase, digits, symbols) has
\[
H = 8 \times \log_2(95) \approx 8 \times 6.57 \approx 52.6 \text{ bits.}
\]
Compare a four-word passphrase such as \texttt{correct horse battery staple} drawn from a
7{,}776-word list (the standard Diceware wordlist, so $R = 7{,}776$ and $L = 4$ words rather than
characters):
\[
H = 4 \times \log_2(7{,}776) \approx 4 \times 12.9 \approx 51.5 \text{ bits.}
\]
The two land at nearly identical entropy despite the passphrase requiring no memorized symbol
substitutions at all. Extending the passphrase to six words pushes it past 77 bits, an increase no
realistic complexity rule on an eight-character password can match. The reason is that $R$ for a
human-chosen ``complex'' password is rarely the full 95-character alphabet in practice: people
overwhelmingly reuse the same handful of substitution patterns (\texttt{a}$\to$\texttt{@},
\texttt{s}$\to$\texttt{\$}). An attacker's cracking dictionary already encodes those patterns,
collapsing the effective $R$ back down toward a much smaller number than the nominal keyboard size
suggests.
\end{example}

\begin{definition}{NIST SP 800-63B}
The 2017 federal digital-identity guideline that reshaped modern password policy: it recommends a
minimum length (8 characters, systems must allow at least 64), checking candidates against known
breach datasets rather than imposing arbitrary complexity rules, and eliminating forced periodic
rotation absent evidence of compromise.\srccite{4.2}{1}
\end{definition}

Forced rotation, in particular, backfires in a specific, well-documented way: users cycle through
predictable families -- \texttt{Summer2024!}, then \texttt{Fall2024!} -- that are barely different
from each other and just as guessable as the original. \textbf{Account lockout} after a threshold
of failures (typically 3--10) closes off online guessing specifically, since an attacker who needs
millions of tries cannot make them interactively. The lockout duration itself is a tuning problem,
though: too short, and scripted slow-guessing simply waits it out; too long or indefinite, and an
attacker can weaponize the lockout itself as a denial-of-service against legitimate users.
\textbf{Password history} rejects a candidate matching the last several previous passwords,
closing off the ``cycle between two passwords'' loophole that a naive rotation policy invites.

Two further attack families exploit scale rather than repeated guessing against one account.
\textbf{Credential stuffing} takes real username/password pairs leaked from an unrelated breach and
replays them against many other sites, betting that some fraction of users reused the same password
everywhere. Screening a new password against a known-breach list (the same NIST-recommended control
described above) blocks this only at the moment a password is chosen; the complementary defense is
watching logins for pairs already known from public breach corpora and forcing a reset on a match.
\textbf{Password spraying} inverts the usual lockout assumption: instead of trying many passwords
against one account, it tries one or two common passwords against thousands of accounts. It stays
under each individual account's lockout threshold by never repeating a guess on the same one.
Per-account lockout does nothing against it -- the defense has to watch the account population as a
whole, rate-limiting failed logins by source or in aggregate rather than per account.

\begin{lstlisting}[language=Python, caption={Validating a candidate password against a policy}]
def validate_password(password, policy, breach_list=frozenset()):
    """policy: dict like {"min_length": int, "history": list[str]}.
       Returns a sorted list of violated rule names."""
    violations = []
    if len(password) < policy["min_length"]:
        violations.append("too-short")
    if password in breach_list:
        violations.append("known-breached-password")
    if password in policy.get("history", []):
        violations.append("password-reuse")
    return sorted(violations)
\end{lstlisting}

Storage matters as much as the policy governing creation. Passwords must never be stored in
plaintext; the correct mechanism is a \textbf{salted cryptographic hash} using a deliberately slow
algorithm -- bcrypt, scrypt, or Argon2. The \textbf{salt}, a unique random value prepended before
hashing, defeats precomputed rainbow tables, and the algorithm's slowness makes each guess
computationally expensive even after a full database theft. MD5 and SHA-1 are unacceptable for
this purpose specifically because they are fast general-purpose hashes -- an attacker with a GPU
can test billions of candidates per second against them.

\begin{example}[The 2012 LinkedIn breach]
In June 2012, an attacker posted roughly 6.5 million password hashes stolen from LinkedIn to a
Russian hacking forum. LinkedIn never officially disclosed the attack vector, though court filings
from the resulting class-action settlement pointed to SQL injection. It later emerged in 2016 that
the same theft had actually covered 117 million accounts.\srccite{4.2}{2} LinkedIn's storage
combined two independent failures at once: the hashes were computed with unsalted SHA-1, and no
deliberately slow key-stretching step (bcrypt, scrypt, or repeated hash iteration) was applied on
top of it. SHA-1 alone is a fast general-purpose hash, so the missing salt let attackers reuse
precomputed rainbow tables across every user at once, and the missing key stretching meant each
individual guess cost a GPU almost nothing to test. Security researchers cracked more than 60\% of
the exposed hashes within a single day of the leak, and the fraction climbed further in the weeks
that followed.\srccite{4.2}{2} Salting alone would have defeated the rainbow-table shortcut; key
stretching alone would have made each guess expensive; LinkedIn had neither, which is why the
breach is the textbook illustration of why both properties must hold simultaneously, not just one.
\end{example}

\begin{misconception}
\textbf{``Changing your password every 90 days improves security.''} Forced rotation produces
predictable families that are trivially related to each other. NIST's current guidance recommends
rotating only when there is actual evidence of compromise, not on a fixed calendar.
\end{misconception}

\section{Multi-Factor Authentication}

A password is a single piece of evidence; once an attacker learns it, through phishing, a breach,
or a guess, nothing else stops them from logging in as the real user.

\begin{definition}{Multi-factor authentication (MFA)}
Requiring proof from two or more independent factor \emph{categories} before granting access:
something you \emph{know} (a password), something you \emph{have} (a phone, a hardware key), and
something you \emph{are} (a fingerprint, a face scan).
\end{definition}

The categories matter more than the count -- entering two different passwords is still ``something
you know'' twice, not true MFA, no matter how much it feels like an extra step. Methods vary sharply
in how resistant they are to interception. \textbf{SMS codes} are better than a password alone, but
the phone \emph{number} is not tightly bound to the phone itself. A \textbf{SIM swap} tricks or
bribes a carrier into porting the victim's number to attacker-controlled hardware, after which every
SMS code arrives on the attacker's device instead. \textbf{TOTP authenticator codes} remove the
SIM-swap risk entirely, since the shared secret lives only on the device. The code can still be
phished and relayed in real time by an adversary-in-the-middle kit, though, because the user is
still capable of typing it into a convincing fake page. \textbf{Push notifications} are convenient
but introduce \textbf{MFA fatigue} (push bombing): an attacker who already holds the password
triggers repeated login attempts until an annoyed or confused user taps ``Approve'' just to make
the notifications stop.

\begin{definition}{FIDO2 / WebAuthn}
An open standard for hardware or platform security keys that cryptographically bind an
authentication response to the exact requesting website domain. This makes the resulting handshake
\textbf{phishing-resistant}: the key simply will not respond to a domain that does not match, no
matter how convincing the fake page looks to a human.
\end{definition}

\begin{example}[Change Healthcare, February 2024: One Portal, No MFA]
Change Healthcare processes a large share of prescription and medical-claims transactions across
the U.S. healthcare system. On February 12, 2024, attackers affiliated with the ALPHV/BlackCat
ransomware operation reached one of its Citrix remote-access portals using a single stolen
username and password -- nothing more.\srccite{4.2}{3} UnitedHealth Group's own CEO later
confirmed to Congress that the portal had no multi-factor authentication enabled at all: a leaked
credential alone was sufficient to reach it. Once inside, the attackers moved laterally for about
nine days before deploying ransomware on February 21. The resulting outage disrupted pharmacy and
claims processing nationwide for weeks. UnitedHealth Group's own disclosures put the direct
response and remediation cost at well over \$2~billion by the end of 2024.\srccite{4.2}{3} The
vulnerability here was not a zero-day or a novel technique. It was the exact gap this section
opened with: a password is a single piece of evidence. Once an attacker has it through phishing, a
breach, or reuse, nothing else stops them from logging in as the real user, unless a second,
independent factor is actually required. One unprotected remote-access portal was enough
to cascade into a nationwide outage. That is why MFA enrollment is treated in this chapter as a
baseline configuration setting, not an optional hardening step reserved for especially sensitive
systems.
\end{example}

Evaluating whether a login satisfies real MFA reduces to two questions. First: do the presented
factors come from two or more genuinely different categories? Second, if so: is at least one of
them resistant to interception, or does the scheme rely on something relayable, SIM-swappable, or
rubber-stampable by an exhausted user?

\begin{lstlisting}[language=Python, caption={Classifying an authentication attempt's MFA strength}]
FACTOR_CATEGORY = {"password": "knowledge", "security_question": "knowledge",
                    "sms_code": "possession", "totp": "possession",
                    "hardware_key": "possession", "fingerprint": "inherence"}
PHISHING_RESISTANT = {"hardware_key"}

def classify_mfa(factors_used):
    """factors_used: list of factor name strings.
       Returns "phishing-resistant-mfa", "mfa", or "single-factor"."""
    categories = {FACTOR_CATEGORY[f] for f in factors_used}
    if len(categories) < 2:
        return "single-factor"
    if PHISHING_RESISTANT & set(factors_used):
        return "phishing-resistant-mfa"
    return "mfa"
\end{lstlisting}

\begin{misconception}
\textbf{``Approving a push notification always means the real user is logging in.''} MFA fatigue
attacks exploit exactly the fact that a user may approve a prompt they never initiated, purely to
stop repeated notifications. Approval is not proof of intent unless the prompt itself shows context
the user actually verifies.
\end{misconception}

\section{Session Security}

Authentication happens once; staying logged in across many requests without retyping a password
every time is what a \textbf{session} exists to solve.

\begin{definition}{Session token (session cookie)}
A long, random, unpredictable string a server issues at successful login, which the browser
resends with every subsequent request; the server treats any request carrying a valid token as
coming from the authenticated user it belongs to.
\end{definition}

That convenience has a cost: the token becomes as valuable as the password itself, since whoever
holds a valid token \emph{is} that user as far as the server can tell. \textbf{Session hijacking}
is the theft and reuse of an already-authenticated token -- through cross-site scripting reading it
directly from the page, through sniffing unencrypted HTTP traffic on a shared network, or through
malware reading browser storage on the victim's own device. It requires no knowledge of the
victim's password at all. \textbf{Session fixation} works differently: instead of stealing an
existing token, the attacker forces the victim to authenticate under a token the attacker already
knows, for instance by sending a link containing a pre-set session identifier. If the server never
changes that ID at login, the attacker's known ID becomes attached to a real authenticated session
the moment the victim logs in.

\begin{example}[Why the defenses for hijacking and fixation differ]
Fixation plants a known ID \emph{before} login; hijacking steals a live ID \emph{after} the victim
is already authenticated. \textbf{Session regeneration} -- issuing a brand-new, freshly random
session ID at the moment of login and again at any privilege change -- defeats fixation completely,
since any ID an attacker planted beforehand is simply discarded and never becomes authenticated. It
does nothing for hijacking, though, which instead requires the token itself to be hard to steal.
The \textbf{Secure} cookie flag restricts transmission to HTTPS; the \textbf{HttpOnly} flag hides
the cookie from JavaScript entirely, so a successful XSS injection still cannot read it; and
\textbf{SameSite} restricts whether the cookie rides along with cross-site requests at all.
\end{example}

\begin{figure}[h]
\centering
\begin{tikzpicture}[
  ev/.style={draw=kickstartermaroon, thick, rounded corners, minimum width=3.3cm, minimum height=0.8cm, align=center, font=\scriptsize},
]
\node[font=\scriptsize\sffamily, kickstartermaroon] at (-2.3,1.1) {Fixation};
\node[ev, fill=kickstartercream] (f1) at (0,1.1) {Attacker sends link\\with known session ID};
\node[ev, fill=kickstartercream] (f2) at (3.6,1.1) {Victim logs in --\\ID never regenerated};
\node[ev, fill=kickstartermaroon!20] (f3) at (7.2,1.1) {Attacker's known ID\\is now authenticated};
\draw[-{Latex[length=2mm]}, kickstartermaroon] (f1) -- (f2);
\draw[-{Latex[length=2mm]}, kickstartermaroon] (f2) -- (f3);
\node[font=\scriptsize\sffamily, kickstartermaroon] at (-2.3,-1.1) {Hijacking};
\node[ev, fill=kickstartercream] (h1) at (0,-1.1) {Victim logs in,\\gets fresh session ID};
\node[ev, fill=kickstartercream] (h2) at (3.6,-1.1) {Attacker steals live ID\\(XSS, sniffing, malware)};
\node[ev, fill=kickstartermaroon!20] (h3) at (7.2,-1.1) {Attacker replays\\the stolen ID};
\draw[-{Latex[length=2mm]}, kickstartermaroon] (h1) -- (h2);
\draw[-{Latex[length=2mm]}, kickstartermaroon] (h2) -- (h3);
\end{tikzpicture}
\caption{Fixation plants a known ID before login; hijacking steals a live ID after login. Session
regeneration defeats fixation specifically; it does nothing against hijacking, which needs the
token itself to be hard to steal.}
\label{fig:4.2-fixation-hijacking}
\end{figure}

\begin{lstlisting}[language=Python, caption={Classifying a session-security incident}]
def classify_session_incident(description):
    """Returns one of "hijacking", "fixation", "improper-timeout",
    "insecure-cookie-flag", or "secure" based on keyword signals."""
    lowered = description.lower()
    if "id set before login" in lowered or "attacker-supplied session id" in lowered:
        return "fixation"
    if "stolen token" in lowered or "xss" in lowered or "sniffed" in lowered:
        return "hijacking"
    if "no timeout" in lowered or "indefinite session" in lowered:
        return "improper-timeout"
    if "missing httponly" in lowered or "missing secure flag" in lowered:
        return "insecure-cookie-flag"
    return "secure"
\end{lstlisting}

Expiration limits how much damage any stolen or fixated token can do, regardless of which theft
path was used. An \textbf{idle timeout} ends a session after a period of inactivity, protecting a
user who walks away from a logged-in device. An \textbf{absolute timeout} ends it after a fixed
total duration regardless of activity, limiting how long a stolen token stays useful even if an
attacker keeps it alive with periodic requests. Logout must be a server-side action for any of this
to matter. A ``logout'' that only deletes the browser's cookie leaves the session record alive on
the server, so a token an attacker captured earlier keeps working indefinitely -- the user's click
on ``logout'' buys no real protection at all.

\begin{misconception}
\textbf{``HTTPS alone makes sessions un-hijackable.''} HTTPS protects the token in transit, but a
cookie missing the HttpOnly flag can still be read directly out of the page's JavaScript context by
a successful XSS injection, with no network sniffing required at all.
\end{misconception}

\section{Least Privilege and Access Control Models}

Once an account is authenticated, a separate question remains: how much should it actually be
allowed to do?

\begin{definition}{Principle of least privilege (PoLP)}
Every subject -- account, process, or service -- should be granted the minimum access required to
perform its function, and nothing more.
\end{definition}

Extra access does not sit harmlessly unused; it expands the \textbf{blast radius} of everything
that can go wrong with that account. A phished credential with read-only access to one shared drive
exposes exactly that drive; the same credential with write access to a production database exposes
the database too, through the identical phishing email. This compounds across three separate
threats that share one root cause. External compromise inherits whatever privileges the stolen
account holds; insider misuse is bounded by whatever access the insider actually has; malware
executing under a user's session inherits that user's privileges outright. An over-privileged
account that opens one infected attachment can hand ransomware write access to systems far beyond
its own workstation.

Organizations implement least privilege through a formal access-control model. \textbf{DAC
(Discretionary Access Control)} lets a resource's owner decide who else can access it -- standard
Unix file permissions are the textbook case. It is flexible but prone to sprawl, since no central
authority enforces consistency as individual owners each grant access independently. \textbf{MAC
(Mandatory Access Control)} moves that decision to a central authority using fixed classification
labels on both subjects and resources, so a user cannot grant access above their own clearance even
if they wanted to. It is common in government and military systems: rigid, but strong against
insider misconfiguration specifically because no individual can override it. \textbf{RBAC
(Role-Based Access Control)} attaches permissions to roles rather than individuals, so a job change
means reassigning a role rather than hand-editing dozens of individual grants -- the most common
model in enterprise IT because it scales. \textbf{ABAC (Attribute-Based Access Control)} decides
access dynamically at request time from attributes of the user, resource, and environment --
``allow only if department is finance \emph{and} the time is business hours \emph{and} the device
is company-managed.'' It is more flexible than RBAC because policy reacts to context, at the cost
of being harder to audit.

\begin{lstlisting}[language=Python, caption={Auditing an access request against least-privilege rules}]
def audit_access_request(request):
    """request: dict like {"has_justification": bool,
       "days_since_last_review": int, "is_time_bounded": bool}.
       Returns "approve", "flag-for-review", or "deny", checked in
       priority order."""
    if not request["has_justification"]:
        return "deny"
    if request["days_since_last_review"] > 180:
        return "flag-for-review"
    if not request["is_time_bounded"]:
        return "flag-for-review"
    return "approve"
\end{lstlisting}

\begin{example}[Privilege creep]
An employee transfers from Sales to Engineering and, as part of onboarding, receives new system
access for the Engineering role. The old Sales access is routinely never revoked, because removing
access requires someone to notice and act, while granting new access is a standard onboarding step
that happens automatically. This silent accumulation -- \textbf{privilege creep} -- undermines
least privilege even in an organization that started out perfectly configured. That is exactly why
a \textbf{periodic access review} exists: on a fixed schedule, every access grant a user holds must
be explicitly re-approved, and anything not re-approved is revoked outright.
\end{example}

Least privilege pairs naturally with \textbf{separation of duties} -- structuring a sensitive
process so no single person controls it end to end, so the person who approves a purchase is never
also the person who can issue the payment. It also pairs with \textbf{just-in-time (JIT) access},
where elevated privilege is requested for a specific, time-bound window rather than sitting active
and exploitable all day, expiring automatically once that window closes.

\begin{misconception}
\textbf{``Giving admin rights to everyone is low-risk if you trust your employees.''} Trust does not
shrink blast radius. A trusted employee's account can still be phished or infected by malware, and
malware inherits the privileges of whoever runs it -- over-provisioning turns any single compromise
into an organization-wide one regardless of how trustworthy the account holder is.
\end{misconception}

\section*{Review Questions}
\begin{reviewquestions}
\item Explain, using search-space math, why a policy that mandates length and unpredictability
outperforms one that mandates complexity (uppercase, digits, symbols) against both brute-force and
dictionary attacks.
\item A user enters a password, then is asked to answer a security question. Explain why this does
not constitute true multi-factor authentication, and identify the minimal change that would make it
qualify.
\item An attacker sends a victim a link containing a pre-set session ID, and the victim then logs in
normally. Explain why this is session fixation rather than session hijacking, and identify the
single server-side control that would have prevented it regardless of what the attacker did.
\item A marketing employee transfers to engineering and receives new access, but their old
marketing-drive access is never revoked. Name this phenomenon and explain the specific organizational
process that is designed to catch and correct it.
\item Explain why RBAC and ABAC are not interchangeable, using the finance-department example from
this chapter to illustrate what ABAC can express that RBAC cannot.
\item In the Change Healthcare breach, a single compromised username and password was sufficient
to reach a remote-access portal and, eventually, cascade into a nationwide outage. Explain
specifically what MFA enrollment would have required of the attacker beyond the stolen password,
and why the absence of that single control is described as the root cause rather than a
contributing factor.
\item Using the fixation-versus-hijacking timeline, explain why deploying only session
regeneration at login, with no HttpOnly or Secure cookie flags, would still leave a system fully
exposed to one of the two attack families covered in this section.
\item A security team notices thousands of failed logins spread across many different accounts,
each attempt trying only one or two common passwords, with no single account ever crossing its
lockout threshold. Name this attack, explain why per-account lockout fails to stop it, and explain
how it differs from credential stuffing, a related attack that also spans many accounts.
\item Explain what storing a salted, slow hash of a password instead of the password itself
protects against when a credential database is stolen. Using the 2012 LinkedIn breach, explain what
it does \emph{not} protect against if the salt or the slow algorithm is missing.
\end{reviewquestions}

\begin{answerkey}
\item Entropy is $H = L \times \log_2(R)$: increasing length $L$ grows the exponent that both
brute-force and dictionary attacks must search. Increasing complexity only grows $R$, the alphabet
size, and human-chosen ``complex'' passwords draw from a far smaller effective $R$ than the nominal
95-character keyboard, because people reuse the same handful of substitution patterns. A long
random passphrase of the same or even shorter apparent length (e.g., four Diceware words at
$H \approx 51.5$ bits versus an eight-character complex password at $H \approx 52.6$ bits) matches or
beats complexity-based policy. Extending the passphrase by two more words adds over 25 bits, an
increase no realistic complexity mandate on a short password can replicate.
\answerlabel{4.2.B}
\item Both factors -- the password and the security question -- fall into the same category,
\emph{something you know}; true MFA requires evidence from two or more independent categories, not
two secrets from the same category. The minimal change is replacing the security question with a
factor from a different category, such as a TOTP code (something you have) or a fingerprint
(something you are).
\answerlabel{4.2.C}
\item This is session fixation because the attacker planted a known session ID \emph{before} the
victim ever authenticated, rather than stealing an already-authenticated token afterward. Session
regeneration -- issuing a brand-new random session ID at the moment of login -- would have
prevented it regardless of the attacker's method. The attacker's pre-planted ID is discarded and
never becomes attached to the authenticated session.
\answerlabel{4.2.D}
\item This phenomenon is \textbf{privilege creep}. A periodic access review is the organizational
process designed to catch it: on a fixed schedule, every access grant a user holds must be explicitly
re-approved, and any access not re-approved -- including the marketing employee's stale drive access
-- is revoked outright.
\answerlabel{4.2.E}
\item RBAC assigns permissions to a role and grants them whenever a user holds that role,
regardless of context. ABAC can express a conditional policy such as ``allow only if department is
finance \emph{and} the time is business hours \emph{and} the device is company-managed'' instead,
which requires evaluating multiple contextual attributes at request time rather than a fixed role
membership. RBAC cannot encode the time-of-day or device-management conditions at all without
creating a combinatorial explosion of narrower roles.
\answerlabel{4.2.E}
\item MFA enrollment would have required the attacker to also possess a second, independent
factor -- a registered device, a hardware key, or a biometric -- something the stolen password
alone could never supply, regardless of how the password was obtained. It is described as the root
cause, not a contributing factor, because every later stage of the breach (lateral movement,
ransomware deployment, the nationwide outage) was contingent on that first successful login
succeeding with a password alone. Had a second factor been required, the identical stolen
credential would have granted the attacker nothing at all, stopping the entire chain at its first
step.
\answerlabel{4.2.C}
\item Session regeneration defeats fixation by discarding any session ID an attacker planted
before login, but it does nothing to protect a token \emph{after} a legitimate login has already
happened. Without HttpOnly, a successful XSS injection can still read the live session cookie
directly out of the page's JavaScript context; without Secure, the cookie can still be transmitted
over an unencrypted connection and sniffed. A system with regeneration alone is therefore still
fully exposed to session hijacking, even though it has completely closed off fixation.
\answerlabel{4.2.D}
\item Password spraying tries one or two common passwords against thousands of accounts, instead of
exhausting many passwords against one account. It deliberately stays under each account's lockout
threshold by never repeating a guess on the same one. Per-account lockout only counts failures on a
single account, so a scheme that never trips that count goes unnoticed; the defense instead has to
watch the account population as a whole, such as rate-limiting total failed logins organization-wide.
Credential stuffing differs in what it guesses with: it replays real username/password pairs leaked
from an unrelated breach, betting on password reuse, rather than guessing common passwords blind --
the matching defense is screening logins against known-breached credential pairs, not lockout.
\answerlabel{4.2.B}
\item Storing a salted, slow hash means a stolen database never hands an attacker plaintext
passwords directly; the attacker only gets something they still have to guess and check, one
candidate at a time, against each hash. It does not protect against the guessing itself. LinkedIn's
2012 breach used unsalted SHA-1 with no slow key-stretching step, so precomputed rainbow tables
cracked the unsalted hashes instantly. A GPU tested billions of guesses per second against the fast
algorithm; more than 60\% of the exposed hashes fell within a single day. Salting alone would
have defeated the rainbow tables; a slow algorithm alone would have made each guess expensive;
LinkedIn had neither, which is why both properties have to hold at once.
\answerlabel{4.2.A}
\end{answerkey}

\section*{Sources}
\begin{enumerate}[label=\lbrack\arabic*\rbrack, leftmargin=2.2em]
\item \srcitem{4.2}{1}National Institute of Standards and Technology, \emph{Special Publication
800-63B: Digital Identity Guidelines -- Authentication and Lifecycle Management}, June 2017.
\item \srcitem{4.2}{2}Brian Krebs, ``As Scope of 2012 Breach Expands, LinkedIn to Again Reset
Passwords for Some Users,'' KrebsOnSecurity, May 18, 2016 (117 million accounts, unsalted,
unstretched SHA-1 hashes); cracking-rate figure per Computerworld, ``Hackers Crack More Than 60\%
of Breached LinkedIn Passwords,'' June 2012, citing Sophos researcher Chester Wisniewski.
\item \srcitem{4.2}{3}Bill Toulas, ``Change Healthcare Hacked Using Stolen Citrix Account With No
MFA,'' BleepingComputer, April 30, 2024; Andrew Witty (CEO, UnitedHealth Group), written testimony
before the U.S.\ House Energy and Commerce Committee, May 1, 2024; cost figure per UnitedHealth
Group's Q2 2024 earnings disclosure, July 16, 2024, updated in subsequent quarterly filings.
\end{enumerate}
